\documentclass[12pt]{article}
\usepackage[utf8]{inputenc}
\usepackage[T1]{fontenc}
\usepackage{hyperref}
\title{Waruna Language Resources for K-12 Curriculum Development}
\author{Compiled by Claude Sonnet 5.5 (claude-sonnet-5-5) for Andr\'as Kornai}
\date{}
\begin{document}
\maketitle

\section{Introduction}
Waruna (ISO 639-3 wrv; Glottocode waru1267) is a Papuan language of Papua New Guinea in the Suki-Gogodala group. It is spoken in a bend of the Fly River, in a single village, Waruna, in Gogodala Rural LLG. Estimates of the speaker population vary. Global Recordings Network lists 600, and Joshua Project gives 1,500 for the Waruna people. Glottolog rates the language as not endangered. The closest relative is Ari, at about 50\% lexical similarity. Very little published material exists, and the searches found no school materials.

\section{Primary Educational Resources}
The searches found no primers, textbooks or school programmes in Waruna. Joshua Project notes a lack of written materials.

\section{Linguistic and Reference Materials}
General reference pages exist, but the searches found no dictionary or grammar.
\begin{itemize}
\item Wikipedia, Waruna language: \url{https://en.wikipedia.org/wiki/Waruna_language}
\item Glottolog entry (waru1267): \url{https://glottolog.org/resource/languoid/id/waru1267}
\item ASJP word list for Waruna: \url{https://asjp.clld.org/languages/WARUNA}
\end{itemize}

\section{Religious and Community Resources}
Joshua Project reports that Bible translation is still needed and that no Bible portions have been reported in Waruna.
\begin{itemize}
\item Joshua Project, Waruna people group: \url{https://joshuaproject.net/people_groups/15877/PP}
\item Joshua Project, language page: \url{https://joshuaproject.net/languages/wrv}
\end{itemize}

\section{Digital and Audio Resources}
Global Recordings Network has a language page for wrv. I did not confirm what recordings it holds.
\begin{itemize}
\item Global Recordings Network: \url{https://globalrecordings.net/en/language/wrv}
\end{itemize}

\section{Curriculum Development Context}
Waruna appears to be a vigorous community language of one village. No orthography, teaching materials or written literature were found. A curriculum would therefore have to start with community-led documentation, orthography development and oral materials. It could also draw on neighbouring Suki-Gogodala languages, though I did not research those.
\end{document}